\documentclass[10pt,a4paper]{amsart}
\usepackage[margin=19mm]{geometry}
\usepackage{amsmath,amssymb,mathtools}
\usepackage[scr=rsfs,cal=euler]{mathalfa}
\usepackage{xfrac}
\usepackage[unicode,hidelinks,bookmarks=false]{hyperref}
\newcommand{\ie}{i.e.\ }
\newcommand{\cf}{cf.\ }
\newcommand{\resp}{respectively\ }
\newcommand{\Subsection}{\subsection}
\providecommand{\nobreakdash}{\nobreak\hbox{-}}
\setlength{\emergencystretch}{2em}
\multlinegap0pt
\usepackage{amssymb}
\usepackage[shortlabels]{enumitem}
\newtheorem{proposition}{Proposition}
\newtheorem{corollary}{Corollary}
\newtheorem{lemma}{Lemma}
\newtheorem{theorem}{Theorem}
\newcommand\C{\mathbb{C}}
\newcommand\E{\mathbb{E}}
\newcommand\GL{{\mathrm{GL}}}
\newcommand\Gfr{\mathscr{G}}
\newcommand\Pbb{\mathbb{P}}
\newcommand\Pfr{\mathscr{P}}
\newcommand\R{\mathbb{R}}
\newcommand\SO{{\mathrm{SO}}}
\newcommand\SU{{\mathrm{SU}}}
\newcommand\Sp{{\mathrm{Sp}}}
\newcommand\Tr{\mathrm{Tr}}
\newcommand\U{{\mathrm U}}
\newcommand\Ufr{\mathscr{U}}
\newcommand\Z{{\mathbb Z}}
\title{The central heat trace on large compact classical groups}
\author{Thibaut Lemoine, Myl\`ene Ma\"ida}
\date{Source: arXiv:2511.08288v2 (CC BY 4.0)}
\begin{document}
\maketitle

\noindent\emph{Source and licence.} The central heat trace on large compact classical groups. Version arXiv:2511.08288v2 (CC BY 4.0), distributed by arXiv under the Creative Commons Attribution 4.0 International licence, http://creativecommons.org/licenses/by/4.0/, as recorded in the arXiv licence metadata for this exact version and shown on the abstract page https://arxiv.org/abs/2511.08288v2. Full licence notice: \texttt{licenses/arxiv-2511.08288-CC-BY-4.0.txt}.\par\smallskip

\noindent\textit{Editorial scope.} Counted unit: the article's theorem thm:asympt\_expansion\_detailed (source0.tex lines 718--735). The article's complete original proof of it is delivered with it (source0.tex lines 739--795, one \texttt{proof} environment of 57 lines, which the article places away from the statement). No mathematics is written, repaired or re-derived: the statement and the proof are the article's own lines, in the article's own notation, and no retained line is substituted or re-flowed.  Retained uncounted as the source-local context the proof needs: lines 370--379; lines 399--404; lines 510--529; lines 534--551; lines 578--593.  Nothing was omitted from any retained span, so the package carries no omission record.  No retained line contains a citation, so the wrapper carries no bibliography and no bibliography entry.  No retained line contains a figure environment, an image-inclusion command or drawing code, so no image is delivered and the package declares no asset dependency.  Editorial formatting only, in addition: an AMS \texttt{amsart} preamble that restates the article's own declarations and macros used by the retained lines, namely the environments \texttt{proposition}, \texttt{corollary}, \texttt{lemma}, \texttt{theorem} on separate counters, together with the standard packages those lines need.  The retained environments print in this document as Proposition 1, Corollary 1, Lemma 1, Theorem 1 in order of appearance, whereas the article prints the same items with the numbering of its own full document.  The complete native source of the article as distributed in the arXiv e-print for this exact version is bundled unchanged as \texttt{sources/188877/source0.tex}.
\begin{proposition}\label{prop:bound_content}
For any $\alpha\in\Pfr$,
\begin{equation}\label{eq:bound_content_1}
\vert K(\alpha)\vert\leq \vert\alpha\vert^2,
\end{equation}
and
\begin{equation}\label{eq:bound_content_2}
-\ell(\alpha)\vert\alpha\vert\leq 2K(\alpha) \leq \alpha_1\vert\alpha\vert.
\end{equation}
\end{proposition}

\begin{proposition}\label{prop:symmetry_partitions}
Let $\alpha,\beta$ two independent $q$-uniform random partitions and $P\in\R[X,Y]$ be a bivariate polynomial. For any integers $k_1,k_2,k_3,k_4\geq 0$, if $k_1$ or $k_2$ is odd, then
\begin{equation}
\E[K(\alpha)^{k_1}K(\beta)^{k_2}P(\vert\alpha\vert,\vert\beta\vert)]=0.
\end{equation}
\end{proposition}

\begin{corollary}\label{cor:heat_trace_expectation}
For all $N$ and all classical groups $G_N$, the central heat trace on $G_N$ can be rewritten as an expectation over random partitions.
\begin{enumerate}
\item If $G_N=\U(N)$ or $\SU(N)$, let $(\alpha,\beta,n)$ be independent random variables such that $\alpha,\beta\sim\Ufr(q_t)$ and $n\sim\Gfr(q_t)$. We have
\begin{align}
\begin{split}
\Tr(e^{\frac{t}{2}\Delta_{\U(N)}}) & =\frac{\theta(q_t)}{\phi(q_t)^2}\E\left[e^{-\frac{t}{N}F^{A'}(\alpha,\beta,n)}\mathbf{1}_{\Lambda_N^{(1)}}(\alpha,\beta,n)\right]\\
\Tr(e^{\frac{t}{2}\Delta_{\SU(N)}}) & = \frac{1}{\phi(q_t)^2}\E\left[e^{-\frac{t}{N}F_1^A(\alpha,\beta)+\frac{t}{2N^2}F_2^A(\alpha,\beta)}\mathbf{1}_{\Lambda_N^{(2)}}(\alpha,\beta)\right].
\end{split}
\end{align}
\item If $G_N=\SO(2N+1),$ $\Sp(N)$ or $\SO(2N)$, let $\mu$ be a single random partition such that $\mu\sim\Ufr(q_t)$. We have
\begin{align}
\begin{split}
\Tr\left(e^{\frac{t}{2}\Delta_{\SO(2N+1)}}\right) & = \frac{1}{\phi(q_t)}\E\left[e^{-\frac{t}{2N+1}F^B(\mu)}\mathbf{1}_{\{\ell(\mu)\leq N\}}\right],\\
\Tr\left(e^{\frac{t}{2}\Delta_{\Sp(N)}}\right) & = \frac{1}{\phi(q_t)}\E\left[e^{-\frac{t}{2N}F^C(\mu)}\mathbf{1}_{\{\ell(\mu)\leq N\}}\right],\\
\Tr(e^{\frac{t}{2}\Delta_{\SO(2N)}}) & =\frac{1}{\phi(q_t)}\E\left[e^{-\frac{t}{2N}F_1^D(\mu)}\sum_{m=0}^\infty\sum_{n=-m}^m e^{-\frac{t}{2}F_{2,N}^D(\mu,m,n)}\mathbf{1}_{\{\ell(\mu)\leq N-2\}}\right].
\end{split}
\end{align}
\end{enumerate}
\end{corollary}

\begin{lemma}\label{lem:lb-casimirs}
Let $G_N$ be a compact classical group of type $\tau\in\{A',A,B,C,D\}$. We have the lower bounds
\[
c_2(\lambda_N^{A'}(\alpha,\beta,n)) \geq \frac12(\vert\alpha\vert+\vert\beta\vert)+\left(n+\frac{\vert\alpha\vert-\vert\beta\vert}{N}\right)^2,\quad \forall (\alpha,\beta,n)\in\Lambda_N^{(1)},
\]
\[
\frac{2}{N}F_1^A(\alpha,\beta)-\frac{1}{N^2}F_2^A(\alpha,\beta)\geq -\frac12(\vert\alpha\vert+\vert\beta\vert),\quad  \forall (\alpha,\beta)\in\Lambda_N^{(2)},
\]
\[
\frac{2}{2N+1}F^B(\mu)\geq -\frac{N+1}{2N+1}\vert\mu\vert \geq -\frac{2}{3}\vert\mu\vert ,\quad \forall \mu\in\Lambda_N^{(3)},
\]
\[
\frac{1}{N}F^C(\mu)\geq-\frac12\vert\mu\vert,\quad \forall\mu\in\Lambda_N^{(3)},
\]
\[
\frac{1}{N}F_1^D(\mu)\geq-\frac12\vert\mu\vert,\quad \forall \mu\in\Pfr\ \text{s.t.}\ \ell(\mu)\leq N-2.
\]
\end{lemma}

\begin{lemma}\label{lem:dev_ineq_bis}
For any integers $n\in\Z$ and $p\geq 1$ and any random $q_t$-uniform $\alpha,\beta$ with $t>0$, we have as $N\to\infty$
\[
\Pbb((\alpha,\beta,n)\notin\Lambda_N^{(1)})=O_t(N^{-p}).
\]
\[
\Pbb((\alpha,\beta)\notin\Lambda_N^{(2)})=O_t(N^{-p}).
\]
\[
\Pbb(\alpha\notin\Lambda_N^{(i)})=O_t(N^{-p}),\quad i\in\{3,4\}.
\]
Furthermore, for any $\gamma\in(0,\infty)$,
\[
\Pbb(\vert\alpha\vert>N^\gamma)=O_t(e^{-\frac{t}{4}N^\gamma}).
\]
\end{lemma}

\begin{theorem}
 \label{thm:asympt_expansion_detailed}
 Let $(G_N)_{N \geq 1}$ be a sequence of compact classical groups of type $\tau \in \{A, A', B, C, D\},$ such that $G_N\subset\GL_N(\C)$. Then $\Tr(e^{\frac{t}{2}\Delta_{G_N}})$ admits the following asymptotic expansion as $N$ tends to infinity~: let $(\alpha, \beta,n)$ be independent random variables such that $\alpha, \beta \sim \Ufr(q_t)$ and $n \sim  \Gfr(q_t),$ we have
 \begin{itemize}
  \item if $G_N = \U(N)$ or $\SU(N),$ ($\tau \in \{A', A\}$ respectively)
  \[ \Tr(e^{\frac{t}{2}\Delta_{G_N}}) = \sum_{k=0}^p \frac{a_{2k}^\tau(t)}{N^{2k}} + O_t(N^{-2p-2}), \]
  with, in the case of $\U(N),$
  \[ a_{2k}^{A'}(t) = \frac{\theta(q_t) t^{2k}}{\phi(q_t)^2 (2k)!} \mathbb E\left[\left(K(\alpha) + K(\beta) +n (|\alpha|-|\beta|)\right)^{2k}\right], \]
  and, in the case of $\SU(N),$
   \[ a_{2k}^A(t) = \frac{1}{\phi(q_t)^2} \sum_{k_1+k_2+k_3 = k} t^{2k-k_3}\frac{\mathbb E\left[K(\alpha)^{2k_1} K(\beta)^{2k_2}(|\alpha|-|\beta|)^{2k_3}\right]}{2^{k_3}(2k_1)!(2k_2)!k_3!}. \]
 \item if $G_N = \SO(N)$ ($\tau \in\{B,D\}$) or $G_N=\Sp(N/2)$ ($\tau=C$),
\[ \Tr(e^{\frac{t}{2}\Delta_{G_N}}) = \sum_{k=0}^p \frac{a_{k}^\tau(t)}{N^{k}} + O_t(N^{-p-1}), \]
with, in the case of $\SO(N)$,
\[  a_{k}^{B}(t) = a_k^D(t)= \frac{(-t)^{k}}{\phi(q_t) k!} \mathbb E\left[\left(K(\alpha) - \frac 1 2 |\alpha|\right)^{k}\right],\]
and, in the case of $\Sp(N/2)$,
\[  a_{k}^{C}(t) = \frac{(-t)^{k}}{\phi(q_t) k!} \mathbb E\left[\left(K(\alpha) + \frac 1 2 |\alpha|\right)^{k}\right].\]
\end{itemize}
\end{theorem}

\begin{proof}[Proof of Theorem~\ref{thm:asympt_expansion_detailed} for $\SU(N)$]\mbox{}\\


\emph{Step 1.} From Corollary~\ref{cor:heat_trace_expectation}, one can write  $\Tr(e^{\frac{t}{2}\Delta_{\SU(N)}})$ as an expectation of an exponential function of $\frac t N$ under the $q_t$-uniform measure:
\[
\Tr(e^{\frac{t}{2}\Delta_{\SU(N)}})=\frac{1}{\phi(q_t)^2}\E[e^{- \frac t N F_{N}(\alpha,\beta)}\mathbf{1}_{\Lambda_N^{(2)}}(\alpha,\beta)], \text{ with } F_{N}(\alpha,\beta)=F_1^A(\alpha,\beta)-\frac{1}{2N}F_2^A(\alpha,\beta).
\]

\emph{Step 2.} We perform a Taylor expansion of the exponential term for fixed $\alpha$ and $\beta$  and show that the remainder is negligible.
The key point will be to use the lower bound on $F_N$ obtained  Lemma \ref{lem:lb-casimirs}. In this case,
\[
e^{- \frac t N F_{N}(\alpha,\beta)}=\sum_{k=0}^{2p+1} \frac{(-t)^k F_{N}(\alpha,\beta)^k}{k!N^k}+ \frac{F_N(\alpha, \beta)^{2p+2}}{N^{2p+2}(2p+1)!} \int_0^t (t-s)^{2p+1} e^{-\frac{s}{N} F_{N}(\alpha,\beta)}ds.
\]
We have to show that $\forall s \leq t,$
\[
I_{N,p}(s):=\mathbb E\left(F_{N}(\alpha,\beta)^{2p+2} e^{-\frac{s}{N} F_{N}(\alpha,\beta)}\mathbf{1}_{\Lambda_N^{(2)}}(\alpha,\beta)\right) =O_t(1),
\]
in the sense that it is bounded by a constant depending on $t$ and $p$ but not depending on $N.$ We have the following upper bounds:
\begin{eqnarray*}
 I_{N,p}(s) & \leq & \mathbb E\left(F_{N}(\alpha,\beta)^{2p+2}\left( 1 +  e^{-\frac{t}{N} F_{N}(\alpha,\beta)}\right)\mathbf{1}_{\Lambda_N^{(2)}}(\alpha,\beta)\right) \\
 & \leq& \mathbb E\left(F_{N}(\alpha,\beta)^{2p+2}\right) +  \mathbb E\left(F_{N}(\alpha,\beta)^{6p+6}\right)^{1/3} \mathbb E \left(  e^{-\frac{3t}{2N} F_{N}(\alpha,\beta)}  \mathbf{1}_{\Lambda_N^{(2)}}(\alpha,\beta) \right)^{2/3}.
\end{eqnarray*}


From Proposition \ref{prop:bound_content}, we have that
\[
 |F_{N}(\alpha,\beta)|  \leq |K(\alpha)| + |K(\beta)| + |\alpha|^2 + |\beta|^2   \leq 2 |\alpha|^2 + 2 |\beta|^2
\]
and it is easy to check that, for any $t >0$ and any $k\in\mathbb N,$ $\sum_{\alpha \in \Pfr} |\alpha|^{k} q_t^{|\alpha|} <\infty$. This leads to a control, depending only on $p,$ on the two quantities $\mathbb E\left(F_{N}(\alpha,\beta)^{2p+2}\right)$ and $\mathbb E\left(F_{N}(\alpha,\beta)^{6p+6}\right)$. Using the lower bound from Lemma \ref{lem:lb-casimirs} which is valid for any $(\alpha, \beta) \in \Lambda_N^{(2)},$ we get
\[   \mathbb E \left(  e^{-\frac{3t}{2N} F_{N}(\alpha,\beta)}  \mathbf{1}_{\Lambda_N^{(2)}}(\alpha,\beta) \right) \leq \mathbb E \left(  e^{\frac{3t}{8}(|\alpha|+|\beta|)}   \right) = \frac{1}{\phi(q_t)^2}\left(\sum_{\alpha \in \Pfr} (q_{t/4})^{|\alpha|}\right)^2.\]
The latter is finite and independent of $N.$

\emph{Step 3.} We write $\mathbf{1}_{\Lambda_N^{(2)}} = 1 - \mathbf{1}_{(\Lambda_N^{(2)})^c}$ in each term of the sum and then use deviation inequalities to discard the terms involving $\mathbf{1}_{(\Lambda_N^{(2)})^c}.$


For each $k \leq2p+1,$  we have
\[
\left\vert\E[F_{N}(\alpha,\beta)^k(1-\mathbf{1}_{\Lambda_N^{(2)}}(\alpha,\beta))]\right\vert\leq \E[F_{N}(\alpha, \beta)^{2k}]^{1/2}   \Pbb((\alpha,\beta)\notin\Lambda_N^{(2)}).
\]
By Lemma~\ref{lem:dev_ineq_bis}, we therefore have that
\[
\E[F_{N}(\alpha,\beta)^k\mathbf{1}_{\Lambda_N^{(2)}}(\alpha,\beta)]=\E[F_{N}(\alpha,\beta)^k]+O_t(N^{-2p-2}).
\]
Altogether,
\[
\Tr(e^{\frac{t}{2}\Delta_{\SU(N)}})=\sum_{k=0}^{2p+1}\frac{ (-t)^k  }{N^k k!\phi(q_t)^2}  \E[F_{N}(\alpha,\beta)^k]+O_t(N^{-2p-2}).
\]
Let us expand the powers of $F_{N}^k$ in order to extract the right powers of $N$ and check that, for any $k\geq 0,$ $\E[F_{N}(\alpha,\beta)^{2k+1}] =0.$ We have
\begin{align*}
\sum_{k=0}^{2p+1}\frac{(-t)^k\E[F_{N}(\alpha,\beta)^k]}{k!N^k}= &  \sum_{k_1+ 2k_2\leq 4p+2} \frac{1}{k_1!k_2!}\frac{(-1)^{k_1}t^{k_1+k_2}}{2^{k_2}N^{k_1+2k_2}}\E[(K(\alpha)+K(\beta))^{k_1}(\vert\alpha\vert-\vert\beta\vert)^{2k_2}]\\
= & \sum_{k_1+2k_2\leq 2p+1}\frac{1}{k_1!k_2!}\frac{(-1)^{k_1}t^{k_1+k_2}}{2^{k_2}N^{k_1+2k_2}}\\
&\times \E[(K(\alpha)+K(\beta))^{k_1}(\vert\alpha\vert-\vert\beta\vert)^{2k_2}] + O_t(N^{-2p-2})\\
= & \sum_{\ell_1+\ell_2+2k_2\leq 2p+1}\frac{(-1)^{\ell_1+\ell_2}t^{\ell_1+\ell_2+k_2}}{\ell_1!\ell_2!k_2!2^{k_2}N^{\ell_1+\ell_2+2k_2}}\\
&\times \E[K(\alpha)^{\ell_1}K(\beta)^{\ell_2}(\vert\alpha\vert-\vert\beta\vert)^{2k_2}]+O_t(N^{-2p-2}).
\end{align*}
In the second equality, we used the fact that the terms such that $2p+2 \leq k_1+2k_2 \leq4p+2$ are absorbed in the $O_t(N^{-2p-2})$, and in the last equality we used the binomial formula. Finally, by Proposition~\ref{prop:symmetry_partitions}, the terms vanish unless both $\ell_1$ and $\ell_2$ are even, and we obtain the statement of the theorem.
\end{proof}

\end{document}
